\documentclass[12pt,a4paper]{article}
\usepackage[utf8]{inputenc}
\usepackage[T2A]{fontenc}
\usepackage[russian]{babel}
\usepackage{amsmath}
\usepackage{graphicx}
\usepackage{geometry}
\usepackage{hyperref}

\geometry{top=2cm,left=2cm,right=2cm,bottom=2cm}
\graphicspath{{../images/}{../graphics/}}
\sloppy

\title{Ответы на контрольные вопросы\\
Лабораторная работа №5: Спектральная теория графов}
\author{Ануфриев Андрей Сергеевич (465029) \\ Бахуров Николай Алексеевич (504862) \\ Раджапов Фёдор Сергеевич (502544)}
\date{01.11.2026}

\begin{document}
\maketitle

\begin{abstract}
Письменные ответы на восемь контрольных вопросов лабораторной работы №5 (Спектральная теория графов), касающиеся графов и их типов, матрицы Лапласа и матрицы смежности, спектральной кластеризации, матрицы инцидентности, векторов вероятности и марковских процессов, а также алгоритма PageRank.
\end{abstract}

\section*{1. Что такое граф? Какие типы графов вы знаете? Что такое компоненты связности?}

Граф $G=(V,E)$ --- математическая структура, состоящая из множества вершин (узлов) $V$ и множества рёбер $E$, соединяющих пары вершин. Он моделирует паровые отношения между объектами.

Основные типы графов:
\begin{itemize}
    \item \emph{Неориентированный граф} --- рёбра не имеют направления; отношение симметрично.
    \item \emph{Ориентированный граф (диграф)} --- рёбра (дуги) имеют направление от одной вершины к другой.
    \item \emph{Взвешенный граф} --- рёбра несут числовые веса.
    \item \emph{Простой граф} --- нет петель и нескольких рёбер между одной парой вершин.
    \item \emph{Мультиграф} --- допускаются несколько рёбер и/или петли.
    \item \emph{Полный граф} --- каждая пара различных вершин соединена ребром.
    \item \emph{Двудольный граф} --- вершины разбиты на две части, рёбра только между частями.
\end{itemize}

\emph{Компоненты связности} --- максимальные подграфы, в которых любые две вершины соединены путём. Граф является связанным, если он имеет ровно одну компоненту связности. В ориентированном графе различают слабую (без учёта направлений дуг) и сильную (с учётом направлений дуг) связность.

\section*{2. Что такое матрица Лапласа, для чего она нужна, как её составить? Что показывают её собственные числа и каковы они?}

(Комбинаторная) матрица Лапласа неориентированного графа --- $L=D-A$, где $A$ --- матрица смежности ($A[i,j]=1$, если вершины $i,j$ смежны, иначе $0$), а $D=\operatorname{diag}(d_1,\dots,d_n)$ --- матрица степеней ($d_i$ = степень вершины $i$). Таким образом, $L[i,i]=d_i$ и $L[i,j]=-1$ для ребра $(i,j)$, иначе $0$.

Матрица Лапласа используется для спектральной кластеризации, разбиения графов, диффузии на сетях и изучения связности. Она симметрична и полуопределена положительно. Её собственные числа удовлетворяют
\[
0=\lambda_1\le\lambda_2\le\dots\le\lambda_n.
\]
Для связанного графа $\lambda_1=0$ --- простое собственное число, а его собственный вектор --- постоянный. Второе собственное число $\lambda_2$ (алгебраическая связность) измеряет, насколько хорошо граф связан; собственные векторы наименьших ненулевых собственных чисел раскрывают кластерную (общественную) структуру --- это основа спектральной кластеризации.

\section*{3. Что такое кластеризация графов? Как найти слабые связи в графе?}

Кластеризация графов --- разбиение вершин на группы (кластеры/сообщества) так, что вершины внутри группы плотно связаны, а связи между группами редки. Это неконтролируемое обучение: метки не заданы заранее.

Слабые связи --- рёбра, удаление которых увеличивает число компонент связности или значительно ослабляет внутрикластерную плотность. Их можно найти следующими способами:
\begin{itemize}
    \item вычислением междустины рёбер (рёбра на многих кратчайших путях --- слабые);
    \item спектральной кластеризацией с использованием вектора Фидлера (собственного вектора $\lambda_2$), чье разделение по знаку указывает на слабо связанные части;
    \item оптимизацией модулярности и алгоритмами обнаружения сообществ.
\end{itemize}

\section*{4. Что такая матрица инцидентности? Как её составить? Как она связана с матрицей Лапласа?}

Для неориентированного графа с $n$ вершинами и $m$ рёбрами (ориентированная) матрица инцидентности $B$ имеет размер $n\times m$: $B[i,e]=1$, если вершина $i$ инцидентна ребру $e$, иначе $0$. Для ориентированного графа $B[i,e]=1$, если ребро $e$ выходит из $i$, $-1$, если оно входит в $i$, иначе $0$.

Связь с матрицей Лапласа:
\[
L = B B^{T}
\]
для ориентированной матрицы инцидентности. Это тождество соединяет структуру инцидентности со спектром матрицы Лапласа.

\section*{5. Что такая матрица смежности? Как её составить? Как она используется для описания марковских процессов?}

Матрица смежности $A$ --- матрица размером $n\times n$ с $A[i,j]=1$ (или количеством ссылок), если есть ребро из $j$ в $i$ (для ориентированного графа) или между $i$ и $j$ (для неориентированного), иначе $0$.

Марковский процесс описывается \emph{матрицей переходов (стахастической матрицей)} $P$, чьи элементы --- вероятности переходов, а столбцы (или строки) суммируются в $1$. Из матрицы смежности получают $P$ путем столбцовой нормировки: $P[i,j]=A[i,j]/\sum_i A[i,j]$. Тогда $P$ --- марковская матрица переходов, а $P^{t}$ дает вероятности переходов за $t$ шагов.

\section*{6. Что такое вектор вероятности? Что такое марковский процесс? Как такие процессы связаны с динамическими системами? Как найти предельный вектор вероятности?}

\emph{Вектор вероятности} --- неотрицательный вектор, элементы которого суммируются в $1$; он описывает распределение по состояниям.

\emph{Марковский процесс} --- стохастический процесс, где будущее состояние зависит только от текущего состояния (отсутствие памяти), а не от истории. Его эволюция управляется матрицей переходов $P$: если $\pi^{(t)}$ --- распределение на шаге $t$, то $\pi^{(t+1)}=P\pi^{(t)}$.

Это линейная дискретная динамическая система на симплексе векторов вероятности. \emph{Предельное (стационарное) распределение} $\pi^*$ удовлетворяет $\pi^*=P\pi^*$, т.е. это собственный вектор $P$ для собственного числа $1$, нормированный в сумме в $1$. Для невырожденной апериодической марковской цепи оно единично и $\pi^{(t)}\to\pi^*$ независимо от начального состояния; его можно найти методом степеней или решением $(P-I)\pi^*=0$ с условием $\sum_i\pi^*_i=1$.

\section*{7. Каковы собственные векторы/собственные числа вероятностной матрицы? Что происходит, когда наибольшее собственное число имеет кратность?}

(Столбцово-)стахастическая/вероятностная матрица имеет все элементы неотрицательными, а каждый столбец суммируется в $1$. По теореме Перрона---Фробениуса, ее наибольшее собственное число по модулю --- $\lambda_{\max}=1$, и существует неотрицательный собственный вектор (стационарное распределение). Все остальные собственные числа удовлетворяют $|\lambda|\le1$.

Если $\lambda=1$ имеет кратность больше $1$ (редуцируемая цепь с несколькими замкнутыми классами общения), стационарное распределение не единственно: каждый замкнутый класс несет свое инвариантное распределение, и предельное состояние зависит от начального состояния (цепь распадается на независимые поглащающие классы).

\section*{8. В чем основная идея PageRank? Как он связан с марковскими процессами? Что происходит, если у вершины нет ссылок? Как с этим бороться?}

PageRank ранжирует веб-страницы по важности: страница важна, если на неё ссылаются важные страницы. Формально, вектор PageRank $r$ --- главный собственный вектор столбцово-стахастической матрицы $M$, построенной по структуре ссылок: $M[i,j]=(\text{ссылки }j\to i)/(\text{исходящих ссылок из }j)$, $Mr=\lambda_{\max}r$, $\lambda_{\max}=1$, $\sum_i r_i=1$. Это стационарное распределение марковской цепи, чье состояние --- текущая страница, а переходы следуют по ссылкам.

PageRank --- марковский процесс: случайный серфер переходит по ссылке с текущей страницы с вероятностью $d$, или перепрыгивает на равномерно случайную страницу с вероятностью $1-d$. Матрица переходов --- $dM+(1-d)\frac{\mathbf{1}\mathbf{1}^{T}}{n}$.

\emph{Проблема висячего узла}: если у страницы нет исходящих ссылок, ее столбец в $M$ состоит из нулей, поэтому $M$ не стахастична и вероятность ``вытекает''; цепь не irreducible и стационарное распределение вырождается.

\emph{Способ борьбы}: заменить нулевой столбец $M$ равномерным вектором $\mathbf{1}/n$ (серфер телепортируется с висячей страницы на любую страницу), или, эквивалентно, использовать формулу с затуханием $dM+(1-d)\frac{\mathbf{1}\mathbf{1}^{T}}{n}$ при $d<1$, которая делает цепь невырожденной и апериодической и гарантирует единичное предельное распределение.

\subsection*{Завершение}
Ответы на все восемь контрольных вопросов приведены выше.

\end{document}
